\clearpage
\section{Data and MC samples}
\label{sec:sample}
We use an $\ee$ collision sample corresponding to an integrated luminosity of 20 $\ifb$, collected by BESIII detector at center-of-mass energy of 3.773 $\gev$, under BOSS-7.1.2 version. The details are listed in Table~\ref{tab:datalumi}, quoted from Ref.~\cite{BESIII_lumi0304, BESIII_lumi_2024}.

\begin{table}[t]
    \begin{center}
        \caption{Beam energies and integrated luminosities of the data samples used in this work, quoted from Ref.~\cite{BESIII_lumi0304, BESIII_lumi_2024}.}
        \vspace{4mm}
        \begin{tabular}{c | c | c | c}
            \hline \hline

            $E_{\text{beam}} (\rm GeV)$            & Integrated luminosity
            ($\text{pb}^{-1}$)                     & Data taking time                               & RunNo                          \\
            \hline
            3.773                                  & \( 2931.8 \pm
            0.2 \pm 13.8 \)~\cite{BESIII_lumi0304} & 2010(round03)                                  & 11414 -- 13988, 14395 -- 14604 \\                                                         &                                                                                    &
            2011(round04)                          & 20448 -- 23454                                                                  \\
            \hline
            3.773                                  & \( 4995.0 \pm 19.0 \) ~\cite{BESIII_lumi_2024} &
            2022(round15)                          & 70522 -- 73929                                                                  \\
            \hline
            3.773                                  & \( 8157.0 \pm 31.0 \) ~\cite{BESIII_lumi_2024} &
            2023(round16)                          & 74031 -- 78536                                                                  \\
            \hline
            3.773                                  & \( 4191.0 \pm 16.0 \) ~\cite{BESIII_lumi_2024} &
            2024(round17)                          & 78615 -- 81094                                                                  \\
            % \hline
            % 4.178                       & \( 3195.0 \pm 32.0 \) [19]     &
            % 2016                        & 43716 -- 45105, 45418 -- 47066                                  \\ 
            \hline \hline
        \end{tabular}
        \label{tab:datalumi}
    \end{center}
\end{table}

Simulated Monte Carlo (MC) samples are employed to study backgrounds, optimize event selection, and estimate detection efficiencies. Two categories of MC samples are utilized in this analysis:

\begin{table}[t]
    \begin{center}
        \caption{Processes and luminosity scales of the inclusive MC samples used in this analysis. Separate samples are available for each data-taking round. Most processes use samples with an integrated luminosity equivalent to ten times that of the data.}
        \vspace{4mm}
        \begin{tabular}{l c}
            \hline \hline
            Process                                         & Luminosity scale \\
            \hline
            $e^+e^- \to \psi(3770) \to D^0\bar{D}^0$        & 10 $\times$      \\
            $e^+e^- \to \psi(3770) \to D^+D^-$              & 10 $\times$      \\
            $e^+e^- \to \tau^+\tau^-$                       & 10 $\times$      \\
            $e^+e^- \to \psi(3770) \to \text{non-}D\bar{D}$ & 10 $\times$      \\
            $e^+e^- \to q \bar{q}$                          & 10 $\times$      \\
            $e^+e^- \to \gamma_{\text{ISR}} \psi(2S)$       & 10 $\times$      \\
            $e^+e^- \to \gamma_{\text{ISR}} J/\psi$         & 10 $\times$      \\
            $e^+e^- \to e^+e^-$                             & 0.25 $\times$    \\
            $e^+e^- \to \gamma \gamma$                      & 3 $\times$       \\
            $e^+e^- \to \mu^+\mu^-$                         & 15 $\times$      \\
            \hline \hline
        \end{tabular}
        \label{tab:inclusiveMC_processes}
    \end{center}
\end{table}

\begin{itemize}
    \item \textbf{The inclusive MC samples} comprehensively model known physics processes at $\sqrt{s}=3.773\gev$, including $D\bar{D}$ production, $\tau^+\tau^-$ pair production, initial-state-radiation (ISR) production of vector charmonium states ($J/\psi$ and $\psi(2S)$), QED processes, and continuum light hadron production. The charm working group provides separate samples for the different data-taking rounds(round03-04, round15, round16, and round17)~\cite{BESIII_MEMO_INCLUSIVE}. In this analysis, we use inclusive MC samples with an integrated luminosity equivalent to ten times that of the data for most processes, except for the QED processes $e^+e^- \to e^+e^-$, $\gamma \gamma$, and $\mu^+\mu^-$, where we use the samples with the scaling factors provided by the charm working group. The processes and the corresponding luminosity scales of the samples used in this analysis are summarized in Table~\ref{tab:inclusiveMC_processes}.

    \item \textbf{Exclusive signal and background MC samples} generated specifically for this analysis using the \textsc{PHOKHARA} event generator~\cite{Czyz2006,Rodrigo2001,Rodrigo2002} (version 9.1). These samples are produced independently for each data-taking round (round03, round04, round15, round16, round17) to account for variations in detector conditions. Sufficient statistics are generated for each process to ensure reliable efficiency determination and background studies. The generated processes include:
          \begin{itemize}
              \item $\ee\to \pi^+\pi^-\pi^0\gamma_{\text{ISR}}$ (signal process)
              \item $\ee\to \pi^+\pi^-\pi^0\pi^0$
              \item $\ee\to \pi^+\pi^-\pi^0\pi^0\gamma_{\text{ISR}}$
          \end{itemize}
          These exclusive samples are used for signal efficiency studies, background characterization, and validation of the analysis methodology.
\end{itemize}